
Our results demonstrate that architectural compatibility is a prerequisite for parameter-efficient fine-tuning, not an outcome of it. Zero-shot evaluation exposes structural constraints before compute investment, preventing wasted experiments on tasks that base model architectures cannot support.

\subsubsection{Key Findings and Interpretation}

The binary success/failure pattern across GLUE tasks---QQP at 38\% (12 percentage points below random), SST-2 at 81\% (31 percentage points above random)---reveals that architectural capabilities are discrete rather than continuous. Mamba-130M either possesses the required information flow mechanism or does not. No intermediate partial capability exists: the model cannot perform bidirectional comparison at all, leading to systematically worse-than-random predictions.

This finding challenges a common assumption in transfer learning: that pretrained models possess general language understanding applicable to any downstream task with appropriate fine-tuning. Our results show that language understanding is not monolithic---it comprises specific capabilities tied to architectural design. Causal state-space models understand language in a unidirectional sense (predicting next tokens, aggregating left-to-right context) but fundamentally lack bidirectional comparison mechanisms. Fine-tuning can specialize existing capabilities but cannot add absent ones.

The SST-2 success (81\% zero-shot accuracy) confirms that checkpoint quality is not the limiting factor. Mamba-130M possesses robust pretrained representations---it simply cannot apply them to architecturally incompatible tasks.

\subsubsection{Implications for PEFT Across Architectures}

Our findings establish a workflow for validating PEFT applicability beyond transformers:

\begin{enumerate}
\item \textbf{Zero-shot evaluation as mandatory gate:} Before investing in fine-tuning infrastructure, test base model on target task without adaptation.
\item \textbf{Interpret below-random accuracy as hard constraint:} If zero-shot accuracy falls below random baseline, flag architectural incompatibility. Fine-tuning will fail.
\item \textbf{Above-random accuracy indicates viability:} Even weak above-baseline performance suggests architectural support. Fine-tuning can specialize these capabilities.
\end{enumerate}

This protocol prevents wasted compute. Training on QQP after observing 38\% zero-shot accuracy would amplify architectural bias rather than overcome it. The model would learn spurious correlations that improve training loss without achieving genuine task competence.

For sub-quadratic architectures, this validation is especially critical. Transformers' bidirectional self-attention supports diverse task types, allowing PEFT research to assume architectural compatibility. Causal state-space models, linear attention, and other efficiency-focused architectures introduce constraints that must be validated per-task.

\subsubsection{Limitations and Scope}

\textbf{Single architecture tested:} We evaluate only Mamba-130M, one instance of causal state-space models. Broader claims require testing additional variants (RWKV, S4, H3). However, the architectural constraint we identify---causal processing prevents symmetric comparison---applies to all autoregressive state-space models, suggesting the failure pattern generalizes.

\textbf{Limited task coverage:} Three GLUE tasks provide sufficient evidence for task-dependent compatibility but do not exhaustively characterize state-space model capabilities. Full GLUE evaluation and additional benchmarks would strengthen claims. Our minimal test suite prioritizes rapid validation over comprehensive benchmarking.

\textbf{Sample size:} 100 examples per task balances statistical power with computational efficiency. QQP and SST-2 results achieve high significance (p $\leq$ 0.003) despite modest sample size. MNLI's marginal result (p = 0.42) suggests larger samples might clarify three-way classification performance, but primary findings are robust.

\textbf{Fine-tuning untested:} We validate the zero-shot gate principle but do not empirically test fine-tuning after failure. Our hypothesis predicts fine-tuning would fail or learn spurious correlations, but confirming this requires additional experiments. The strong prior evidence (38\% zero-shot accuracy) makes failure highly likely, but empirical validation remains future work.

\textbf{Generalization beyond Mamba:} While we focus on causal state-space models, other sub-quadratic architectures may exhibit different failure modes. RWKV's bidirectional variants might succeed on bidirectional tasks. Our claims are bounded to causal state-space models; broader architectural families require separate validation.

\subsubsection{Broader Impact}

Our work reduces computational waste by identifying failed experiments early. This efficiency benefit has environmental implications---preventing unnecessary training runs conserves energy and reduces carbon emissions. Practitioners can validate architectural compatibility in minutes rather than discovering failures after hours of fine-tuning.

We do not anticipate negative societal impacts. The zero-shot gate helps practitioners avoid failed experiments but does not enable harmful applications. Preventing wasted compute makes PEFT more accessible to resource-constrained researchers.

However, practitioners might misinterpret our findings as ''state-space models are incompetent'' rather than ''causal state-space models and bidirectional tasks are incompatible.'' Architectural constraints are task-specific, not absolute limitations. Mamba-130M excels at generation-aligned tasks (SST-2: 81\%) and long-context modeling. Our work characterizes compatibility boundaries, not model inferiority.

\subsubsection{Future Work}

\textbf{Bidirectional state-space model validation:} Testing bidirectional state-space models (e.g., RWKV with dual-direction mode) on GLUE tasks would validate whether dual-direction mechanisms overcome limitations we identified while preserving efficiency benefits.

\textbf{Task-architecture taxonomy:} Developing systematic mapping between task requirements (bidirectional comparison, long-range dependencies, structural reasoning) and architectural capabilities (attention mechanisms, recurrence patterns, positional encodings) would guide architecture selection before experimentation.

\textbf{Architectural augmentation:} Targeted modifications might enable adaptation on incompatible tasks. Adding small bidirectional attention layers to causal state-space models might enable paraphrase detection while preserving efficiency.

\textbf{Zero-shot calibration:} Our binary threshold (above/below random baseline) is crude. Future work could develop more nuanced calibration: how far above baseline must zero-shot accuracy be to predict successful fine-tuning? Marginal above-baseline performance may indicate fragile compatibility requiring careful hyperparameter tuning.

\textbf{Broader architecture families:} Extending analysis to linear attention, RWKV, and hybrid models would validate whether the zero-shot gate generalizes beyond causal state-space models. Each architecture class may exhibit unique failure modes requiring tailored compatibility checks.

The zero-shot architectural compatibility gate represents a paradigm shift for PEFT: from ''fine-tune and hope'' to ''validate then fine-tune.'' As efficient architectures proliferate, this principled workflow prevents wasted compute while guiding practitioners toward compatible model-task pairings.
